\section{Dataset \& Tasks}
\label{sec:dataset}

\subsection{Source Data}
To leverage Robo2VLM-1~\cite{chen2025robo2vlm} keyframe selection of the scene diverse DROID~\cite{khazatsky2024droid} dataset.
\cite{chen2025robo2vlm} use a novel framework to identify keyframes within each
trajectory and pair them with the goal text, ordered action-phase labels, and
sensor-derived success/failure signals captured from three synchronized camera views
(top-left, top-right, wrist). From this pool we select 47 scenes from the DROID subset,
covering diverse manipulation tasks such as placing, screwing, folding, and stacking. We label and sort the action phases manually to ensure (1) high quality and diverse scene selection (2) scene and task understanding is humanly possible (3) securing groundtruth annotation on which view holds necesssary visual signal (4) temporal consistency.
As a result we obtain an ordered list of keyframes, their action-phase labels of what the robot is doing, as well as annotation of which view enables comprehension and correct answering. 

We define four multiple-choice failure-mode sub-tasks, all evaluated on real DROID trajectory keyframes.

\textbf{Action Phase Identification} (\texttt{action\_phase\_id}): given a single moment, the goal text, and a claimed current action phase, the model selects which phase (A--D) is currently depicted, or answers ``Cannot be determined'' (E). Random baseline: 33.3\%. 

\textbf{Phase Success} (\texttt{phase\_success}): given a single moment, the goal text, and
a target phase, the model judges whether that phase is Before, Is currently performing,
After, or Cannot be determined. Random baseline: 25.0\%. 

\textbf{Progress Detection} (\texttt{progress}): given two distinct moments and the
goal text, the model judges whether the task progressed forward, backward, or
Cannot be determined. Random baseline: 33.3\%.

\textbf{Task Success} (\texttt{task\_success}): given a single image and the goal text,
the model answers Yes (complete), No (failed or incomplete), or Cannot be determined.
Random baseline: 33.3\%.

For \texttt{action\_phase\_id}, \texttt{task\_success}, and \texttt{progress} we
additionally inject an adversarial \emph{Cannot be determined} (CnBD) variant: when
\texttt{special\_image == "random\_scene"}, the input image is replaced with an unrelated
scene, making ``Cannot be determined'' the only correct answer. This variant tests
whether a model detects a domain mismatch between image and goal text, rather than
guessing a plausible-looking phase.
Additionally we provide each question with context of a sorted list of action phase labels of the scene, to provide the model with potentially needed context. E.g. for action phase detection, the model might answer incorrectly, because it would attribute it sligthley differently formulated action phase.


% Figure: evaluation target examples
\begin{figure}[t]
  \centering
  \includegraphics[width=\linewidth]{../poster_3_D/phase_success_top.pdf}
  \caption{Example questions for the four failure-mode detection sub-tasks.
  \TODO{Write caption describing what each row shows.}}
  \label{fig:task_examples}
\end{figure}

\subsection{Failure-Mode Detection Sub-Tasks}
For failure detection, we prompt the model with all 3 views merged into a single image, as prodvided by \cite{chen2025robo2vlm}. CoT shows the models can identify the individual views in a merged image.
\subsection{Multi-View Consistency}
The same question is additionally posed from multiple camera views of the same scene, if a human labelled the view as "action-phase is comprehensible from single view":
top-left (TL), top-right (TR), and wrist (W). 
\begin{figure}[t]
  \centering
  \includegraphics[width=\linewidth]{../poster_3_D/mutlview_consistency_evaluation_targets_bottom.pdf}
  \caption{Multi-view consistency evaluation: same scene captured from top-left, top-right,
  and wrist cameras. \TODO{Write caption.}}
  \label{fig:multiview_examples}
\end{figure}

\subsection{Dataset Statistics}

\begin{table}[h]
  \centering
  \caption{Dataset statistics. The LoRA split uses a scene-level partition with no
  scene overlap between train, val, and test.}
  \label{tab:dataset_stats}
  \begin{tabular}{lrrr}
    \toprule
    & \textbf{Train} & \textbf{Val} & \textbf{Test} \\
    \midrule
    Scenes      & 20    & 7   & 20   \\
    Questions   & 2535  & 983 & 2706 \\
    \midrule
    \multicolumn{4}{l}{\small Total: 47 scenes, 6224 questions} \\
    \bottomrule
  \end{tabular}
\end{table}

Table~\ref{tab:dataset_stats} reports scene and question counts for the LoRA train/val/test
split. The full benchmark, used for zero-shot and CoT evaluation, totals 47 scenes and
6224 questions; the LoRA split is scene-disjoint, so no scene appears in more than one of
train, val, and test.
% TODO: Add per-sub-task question count breakdown if available.
